% GENERATED FILE. Edit canonical lessons, then run npm run generate:books.
% canonical-source-hash: fnv1a64:ad8389019a024ae9
% canonical-lessons: RU-C458-ravnina, RU-C458-dolina, RU-C458-sklon, RU-C458-vershina, RU-C458-ushchelye, RU-R458-review-451-458

\chapter{Mountains and valleys}
\label{ch:ru-mountains-and-valleys}
\hlchaptermodality{458}

\begin{chapteropening}
\textbf{By the end of this chapter:} \emph{I can describe the land: a plain, a valley, a slope, a summit and a gorge.}

Complete the last lesson of chapter 458: I can describe the land: a plain, a valley, a slope, a summit and a gorge.
\end{chapteropening}

\section[\texorpdfstring{\ru{равнина} (ravnína)}{равнина (ravnína)}]{\ru{равнина} (ravnína) — a plain (flat land)}

\label{lesson:RU-C458-ravnina}



\begin{quote}
Before the new one: say the Russian for to deliver, then the Russian for a courier.
\end{quote}

\subsection*{You'll want to know: \ru{равнина}}
\textbf{\ru{равнина}} (\emph{ravnína}) — \textquotedblleft{}a plain (flat land)\textquotedblright{}.

\begin{quote}
\ru{За} \ru{рекой} \ru{начинается} \ru{широкая} \ru{равнина}.
\end{quote}

\begin{quote}
\emph{Za rekóy nachináyetsya shirókaya ravnína.}
\end{quote}

\begin{quote}
Beyond the river a wide plain begins.
\end{quote}

Feminine. \emph{\ru{Равнина}} has the Church Slavonic \emph{\ru{ра}-} where native \emph{\ru{ровный}} has \emph{\ru{ро}-}.

\begin{cousinweb}
\emph{\ru{равнина}} is built on \emph{\ru{равный}}, the Church Slavonic shape of native \emph{\ru{ровный}} 'level, even', with \emph{-\ru{ина}}. EtymDB takes \emph{\ru{ровный}} from Proto-Slavic \emph{orvĭnŭ}.
\end{cousinweb}

\begin{grammarlens}[title={What you've built}]
One more word for telling, explaining and describing at B1.
\end{grammarlens}

\subsection*{Your turn}
\begin{itemize}
\raggedright
  \item \emph{Say it:} \emph{ravnína}
  \item \emph{Say it:} \emph{ravnína}, once more
  \item \emph{Say it:} a sentence of your own with \emph{ravnína}
  \item \emph{Recall:} say the Russian for to deliver, then the Russian for a courier, then say \emph{ravnína} again
\end{itemize}

\begin{tcolorbox}[breakable,colback=teal!4,colframe=teal!35!black,title={Before you move on}]
What does \ru{равнина} mean? (\textquotedblleft{}A plain (flat land)\textquotedblright{}.) Say it once more.
\end{tcolorbox}
\section[\texorpdfstring{\ru{долина} (dolína)}{долина (dolína)}]{\ru{долина} (dolína) — a valley}

\label{lesson:RU-C458-dolina}



\begin{quote}
Before the new one: say the Russian for a courier, then the Russian for a plain.
\end{quote}

\subsection*{You'll want to know: \ru{долина}}
\textbf{\ru{долина}} (\emph{dolína}) — \textquotedblleft{}a valley\textquotedblright{}.

\begin{quote}
\ru{В} \ru{долине} \ru{течёт} \ru{маленькая} \ru{река}.
\end{quote}

\begin{quote}
\emph{V dolíne techyót málenkaya reká.}
\end{quote}

\begin{quote}
A small river flows through the valley.
\end{quote}

Feminine.

\begin{cousinweb}
\emph{\ru{долина}} is built on the old word \emph{\ru{дол}} 'a valley', still heard in \emph{\ru{по} \ru{горам}, \ru{по} \ru{долам}} 'over hill and dale'. \emph{\ru{Дол}} is inherited from Proto-Slavic \emph{dolŭ} 'a dale, a valley' (EtymDB); EtymDB lists \emph{\ru{долина}} itself without a source.
\end{cousinweb}

\begin{grammarlens}[title={What you've built}]
One more word for telling, explaining and describing at B1.
\end{grammarlens}

\subsection*{Your turn}
\begin{itemize}
\raggedright
  \item \emph{Say it:} \emph{dolína}
  \item \emph{Say it:} \emph{dolína}, once more
  \item \emph{Say it:} a sentence of your own with \emph{dolína}
  \item \emph{Recall:} say the Russian for a courier, then the Russian for a plain, then say \emph{dolína} again
\end{itemize}

\begin{tcolorbox}[breakable,colback=teal!4,colframe=teal!35!black,title={Before you move on}]
What does \ru{долина} mean? (\textquotedblleft{}A valley\textquotedblright{}.) Say it once more.
\end{tcolorbox}
\section[\texorpdfstring{\ru{склон} (sklon)}{склон (sklon)}]{\ru{склон} (sklon) — a slope}

\label{lesson:RU-C458-sklon}



\begin{quote}
Before the new one: say the Russian for a plain, then the Russian for a valley.
\end{quote}

\subsection*{You'll want to know: \ru{склон}}
\textbf{\ru{склон}} (\emph{sklon}) — \textquotedblleft{}a slope\textquotedblright{}.

\begin{quote}
\ru{Зимой} \ru{на} \ru{этом} \ru{склоне} \ru{катаются} \ru{на} \ru{лыжах}.
\end{quote}

\begin{quote}
\emph{Zimóy na étom sklóne katáyutsya na lýzhakh.}
\end{quote}

\begin{quote}
In winter people ski on this slope.
\end{quote}

Masculine. \emph{\ru{На} \ru{склоне} \ru{лет}} — in old age.

\begin{cousinweb}
\emph{\ru{склон}} is built on \emph{\ru{склонить}} 'to bend, to incline', from \emph{\ru{клонить}} 'to lean'. EtymDB lists \emph{\ru{склон}} without a deeper source.
\end{cousinweb}

\begin{grammarlens}[title={What you've built}]
One more word for telling, explaining and describing at B1.
\end{grammarlens}

\subsection*{Your turn}
\begin{itemize}
\raggedright
  \item \emph{Say it:} \emph{sklon}
  \item \emph{Say it:} \emph{sklon}, once more
  \item \emph{Say it:} a sentence of your own with \emph{sklon}
  \item \emph{Recall:} say the Russian for a plain, then the Russian for a valley, then say \emph{sklon} again
\end{itemize}

\begin{tcolorbox}[breakable,colback=teal!4,colframe=teal!35!black,title={Before you move on}]
What does \ru{склон} mean? (\textquotedblleft{}A slope\textquotedblright{}.) Say it once more.
\end{tcolorbox}
\section[\texorpdfstring{\ru{вершина} (vershína)}{вершина (vershína)}]{\ru{вершина} (vershína) — a summit, a peak}

\label{lesson:RU-C458-vershina}



\begin{quote}
Before the new one: say the Russian for a valley, then the Russian for a slope.
\end{quote}

\subsection*{You'll want to know: \ru{вершина}}
\textbf{\ru{вершина}} (\emph{vershína}) — \textquotedblleft{}a summit, a peak\textquotedblright{}.

\begin{quote}
\ru{До} \ru{вершины} \ru{горы} \ru{мы} \ru{шли} \ru{пять} \ru{часов}.
\end{quote}

\begin{quote}
\emph{Do vershíny gorý my shli pyat chasóv.}
\end{quote}

\begin{quote}
It took us five hours to walk to the top of the mountain.
\end{quote}

Feminine. Built on \emph{\ru{верх}} — the top.

\begin{cousinweb}
\emph{\ru{вершина}} is built on \emph{\ru{верх}} 'the top' with \emph{-\ru{ина}}. EtymDB lists \emph{\ru{вершина}} without a deeper source.
\end{cousinweb}

\begin{grammarlens}[title={What you've built}]
One more word for telling, explaining and describing at B1.
\end{grammarlens}

\subsection*{Your turn}
\begin{itemize}
\raggedright
  \item \emph{Say it:} \emph{vershína}
  \item \emph{Say it:} \emph{vershína}, once more
  \item \emph{Say it:} a sentence of your own with \emph{vershína}
  \item \emph{Recall:} say the Russian for a valley, then the Russian for a slope, then say \emph{vershína} again
\end{itemize}

\begin{tcolorbox}[breakable,colback=teal!4,colframe=teal!35!black,title={Before you move on}]
What does \ru{вершина} mean? (\textquotedblleft{}A summit, a peak\textquotedblright{}.) Say it once more.
\end{tcolorbox}
\section[\texorpdfstring{\ru{ущелье} (ushchélye)}{ущелье (ushchélye)}]{\ru{ущелье} (ushchélye) — a gorge, a ravine}

\label{lesson:RU-C458-ushchelye}



\begin{quote}
Before the new one: say the Russian for a slope, then the Russian for a summit.
\end{quote}

\subsection*{You'll want to know: \ru{ущелье}}
\textbf{\ru{ущелье}} (\emph{ushchélye}) — \textquotedblleft{}a gorge, a ravine\textquotedblright{}.

\begin{quote}
\ru{Дорога} \ru{идёт} \ru{по} \ru{узкому} \ru{ущелью}.
\end{quote}

\begin{quote}
\emph{Doróga idyót po úzkomu ushchélyu.}
\end{quote}

\begin{quote}
The road runs through a narrow gorge.
\end{quote}

Neuter.

\begin{cousinweb}
\emph{\ru{ущелье}} is usually traced to \emph{\ru{щель}} 'a crack, a gap', with \emph{\ru{у}-} and \emph{-\ru{ье}}: a deep crack between mountains. EtymDB does not list it, so the line is hedged.
\end{cousinweb}

\begin{grammarlens}[title={What you've built}]
One more word for telling, explaining and describing at B1.
\end{grammarlens}

\subsection*{Your turn}
\begin{itemize}
\raggedright
  \item \emph{Say it:} \emph{ushchélye}
  \item \emph{Say it:} \emph{ushchélye}, once more
  \item \emph{Say it:} a sentence of your own with \emph{ushchélye}
  \item \emph{Recall:} say the Russian for a slope, then the Russian for a summit, then say \emph{ushchélye} again
\end{itemize}

\begin{tcolorbox}[breakable,colback=teal!4,colframe=teal!35!black,title={Before you move on}]
What does \ru{ущелье} mean? (\textquotedblleft{}A gorge, a ravine\textquotedblright{}.) Say it once more.
\end{tcolorbox}
\section[(dialogue)]{Words from chapters 451–458 — a review}

\label{lesson:RU-R458-review-451-458}



\begin{quote}
Say the Russian for an embassy and a passport for travel abroad.
\end{quote}

\subsection*{Your turn}
Say these aloud:

\begin{itemize}
\raggedright
  \item an embassy, a passport for travel abroad, a small parcel, an addressee, to issue
  \item a bakery, a grocery store, a counter, a shop assistant, to weigh
  \item a fitting room, to try on, a size, to pinch, to exchange
  \item a sale, a markdown, a special offer, to haggle, to overpay
  \item to settle up, payment, a bank card, to break, to shortchange
  \item a bureau de change, a foreign currency, a warranty, faulty, to save
  \item to pack, to wrap, delivery, to deliver, a courier
  \item a plain, a valley, a slope, a summit, a gorge
\end{itemize}

\begin{tcolorbox}[breakable,colback=teal!4,colframe=teal!35!black,title={Before you move on}]
An embassy? (\textbf{\ru{посольство}}.) A passport for travel abroad? (\textbf{\ru{загранпаспорт}}.) A small parcel? (\textbf{\ru{бандероль}}.) An addressee? (\textbf{\ru{получатель}}.) To issue? (\textbf{\ru{выдавать}}.) A bakery? (\textbf{\ru{булочная}}.) A grocery store? (\textbf{\ru{гастроном}}.) A counter? (\textbf{\ru{прилавок}}.) A shop assistant? (\textbf{\ru{продавщица}}.) To weigh? (\textbf{\ru{взвешивать}}.) A fitting room? (\textbf{\ru{примерочная}}.) To try on? (\textbf{\ru{примерять}}.) A size? (\textbf{\ru{размер}}.) To pinch? (\textbf{\ru{жать}}.) To exchange? (\textbf{\ru{обменивать}}.) A sale? (\textbf{\ru{распродажа}}.) A markdown? (\textbf{\ru{уценка}}.) A special offer? (\textbf{\ru{акция}}.) To haggle? (\textbf{\ru{торговаться}}.) To overpay? (\textbf{\ru{переплачивать}}.) To settle up? (\textbf{\ru{расплачиваться}}.) Payment? (\textbf{\ru{оплата}}.) A bank card? (\textbf{\ru{карточка}}.) To break? (\textbf{\ru{разменивать}}.) To shortchange? (\textbf{\ru{обсчитывать}}.) A bureau de change? (\textbf{\ru{обменник}}.) A foreign currency? (\textbf{\ru{валюта}}.) A warranty? (\textbf{\ru{гарантия}}.) Faulty? (\textbf{\ru{бракованный}}.) To save? (\textbf{\ru{экономить}}.) To pack? (\textbf{\ru{упаковывать}}.) To wrap? (\textbf{\ru{заворачивать}}.) Delivery? (\textbf{\ru{доставка}}.) To deliver? (\textbf{\ru{доставлять}}.) A courier? (\textbf{\ru{курьер}}.) A plain? (\textbf{\ru{равнина}}.) A valley? (\textbf{\ru{долина}}.) A slope? (\textbf{\ru{склон}}.) A summit? (\textbf{\ru{вершина}}.) A gorge? (\textbf{\ru{ущелье}}.)
\end{tcolorbox}
